\documentclass[12pt]{article}
\usepackage{graphicx}
\usepackage[margin=1in]{geometry}
\usepackage{caption}
\usepackage{xcolor}
\usepackage{float}
\usepackage{amsmath}
\usepackage{array}
\usepackage{multicol}
\usepackage{hyperref}
\usepackage{url}
\usepackage{listings}
\usepackage{courier}

\lstset{
  basicstyle=\ttfamily\footnotesize,
  backgroundcolor=\color{gray!10},
  frame=single,
  breaklines=true,
  captionpos=b
}

\setlength{\columnsep}{1cm}

\begin{document}

% Title Block
\begin{figure}[H]
    \begin{minipage}{0.45\textwidth}
        \includegraphics[width=\textwidth]{iiitb_comet.jpeg} % Replace with your image
    \end{minipage} \hfill
    \begin{minipage}{0.45\textwidth}
        \textbf{Name: S. Rithrisha Reddy} \\
        \textbf{Batch: COMETFWC028} \\
        \textbf{Date: 27 May 2025}
    \end{minipage}
\end{figure}

\begin{center}
    {\LARGE \textbf{\textcolor{cyan}{GATE 2009, ECE Question Number 36}}}
\end{center}

\begin{multicols}{2}

\vspace{1em}
\noindent\textbf{Abstract} \\
This project implements the logic equation given in a GATE 2009 ECE question using AVR-GCC and physical hardware (ATmega328P/Arduino). Logical analysis determines the value of variable Z, and the expression is validated using an LED output.

\vspace{1em}
\noindent\textbf{1. Question and Simplification}
\begin{itemize}
    \item Given: \( X = 1 \)
    \item Expression: \( (X + Z)(Y + (Z + Y))(X + Z(X + Y)) = 1 \)
    \item Simplifies to: \( (1)(Y + Z)(1 + Z(1 + Y)) = 1 \)
    \item Solving: \( Z(1 + Y) = 0 \Rightarrow Z = 0 \)
\end{itemize}
\textbf{Correct Answer: (D) Z = 0}

\vspace{1em}
\noindent\textbf{2. Components}
\begin{table}[H]
\small
\centering
\begin{tabular}{|p{4.2cm}|c|}
\hline
\textbf{Component} & \textbf{Qty} \\
\hline
ATmega328P / Arduino UNO & 1 \\
Push Buttons (X, Y, Z) & 3 \\
LED for Output & 1 \\
Resistors (220$\Omega$ for LED, 10k$\Omega$ for pull-downs) & 4 \\
Breadboard & 1 \\
Jumper Wires & 10 \\
Laptop with AVR-GCC & 1 \\
\hline
\end{tabular}
\caption*{Table 1: List of components used}
\end{table}

\vspace{1em}
\noindent\textbf{3. Setup and Connections}
\begin{itemize}
    \item Connect buttons to PD0, PD1, PD2 for inputs X, Y, Z.
    \item Connect LED to PB0 with 220$\Omega$ resistor.
    \item All buttons should use 10k$\Omega$ pull-down resistors.
    \item Common GND and VCC lines shared across all components.
\end{itemize}

\vspace{1em}
\noindent\textbf{4. Logic Summary}
\begin{itemize}
    \item Inputs: X, Y, Z (button-controlled)
    \item Expression: \( (X + Z)(Y + (Z + Y))(X + Z(X + Y)) \)
    \item With X = 1, simplify to: \( (Y + Z)(1 + Z(1 + Y)) \)
    \item Final condition: LED ON if entire expression evaluates to 1
\end{itemize}

\vspace{1em}
\noindent\textbf{5. Pin Mapping}
\begin{itemize}
    \item \textbf{X} – PD0 (Input)
    \item \textbf{Y} – PD1 (Input)
    \item \textbf{Z} – PD2 (Input)
    \item \textbf{LED} – PB0 (Output)
\end{itemize}

\end{multicols}

\newpage
\section*{6. Hardware Code – AVR GCC (C)}
\begin{lstlisting}[language=C]
#define F_CPU 1000000UL
#include <avr/io.h>
#include <util/delay.h>

int main(void)
{
    DDRD = 0x00; // Inputs: PD0 - X, PD1 - Y, PD2 - Z
    DDRB = 0xFF; // Output: PB0 - LED
    PORTD = 0x00;

    while (1)
    {
        unsigned char X = (PIND & (1 << PD0)) >> PD0;
        unsigned char Y = (PIND & (1 << PD1)) >> PD1;
        unsigned char Z = (PIND & (1 << PD2)) >> PD2;

        unsigned char part1 = (X || Z);
        unsigned char part2 = (Y || (Z || Y));
        unsigned char part3 = (X || (Z && (X || Y)));

        unsigned char output = part1 && part2 && part3;

        if (output)
            PORTB |= (1 << PB0); // LED ON
        else
            PORTB &= ~(1 << PB0); // LED OFF

        _delay_ms(200);
    }
    return 0;
}
\end{lstlisting}

\section*{7. Analysis}

\subsection*{7.1 Truth Table (sample)}
\begin{table}[H]
\centering
\begin{tabular}{|c|c|c|c|}
\hline
X & Y & Z & Output \\
\hline
1 & 0 & 0 & 1 \\
1 & 1 & 0 & 1 \\
1 & 0 & 1 & 0 \\
1 & 1 & 1 & 0 \\
\hline
\end{tabular}
\caption*{Table 2: Logic validation for Z}
\end{table}

\subsection*{7.2 Circuit Diagram}
\begin{figure}[H]
\centering
\includegraphics[width=0.8\linewidth]{output1.jpeg}
\caption*{Figure: Practical Circuit Using AVR}
\end{figure}

\section*{8. Conclusion}
This project demonstrated evaluation of a GATE logic question both theoretically and practically using AVR-GCC. The LED output confirms that the expression is valid only when \( Z = 0 \), verifying the analytical result.

\end{document}
